\chapter{The Model}
\label{sec:model}

\section{Cells, paths and mode crossing}
\label{sec:model-cells}

A beni value at run time is a JavaScript value. The ones that matter here are the arrays that back
lists. A \emph{cell} $c$ is one such array; two list values that are the same cell are \code{===},
and editing a cell in place changes what every reference to it sees. A list value is a cell or a
\emph{view} of one, a header \code{\{ b, o, length \}} over a base array $b$
(\S\ref{sec:bg-lists}). A view is a holder of its base, and writing through a view writes the base
(\S\ref{sec:hard-views}). A closure that may be called only once is also a cell, for reasons given in
\S\ref{sec:model-functions}: calling it is the write.

Every other value either contains cells or does not. A \emph{path} is a chain of steps from a value
into it:
\begin{itemize}
\item a record field $.f$, a tuple index $.i$, a constructor argument $C\#i$;
\item a list position $[\ast]$, which stands for every element;
\item the \emph{environment} step $\env$, from a function value to everything its closure captured;
\item an \emph{opaque} step $\langle T\rangle$, from a value of a type $T$ whose constructors another
module hides, to every cell inside it.
\end{itemize}
Paths are those of the write-set analysis (\S\ref{sec:bg-browser}) extended by the last two steps.
They are bounded in two ways. A path that re-enters a type already on it is \emph{folded}: the steps
since the first visit are dropped, so every list held anywhere in a tree of type
\code{type Tree = Node String (List Tree)} is reached by one path, and the number of paths under a type
is at most the number of (type, step) edges in the graph of its definition---a recursive type
contributes its definition once, not once per level. Paths of a non-recursive type are additionally
cut at $k$ steps ($k = 8$ in our design), a value deeper than the cut being approximated by $\top$
below. A path is \emph{tracked} when the checker's solved type at its end is \code{List a} or a
function type; only tracked paths carry cells.

Folding a recursive type loses the distinction between levels in the same way that $[\ast]$ loses the
distinction between positions, and the exclusive-contents flag of \S\ref{sec:model-excl} recovers it
for the same edits.

\paragraph{Mode crossing.} A value whose type can contain no cell---\code{Int}, \code{String}, a record
of them, \code{Order}---is never tracked. It \emph{crosses} the ownership axis in the sense of Peters
et al., for whom a type crosses an axis ``when all of its values have no interaction with that axis''
\citep[\S2.3]{peters2026mode}; in their system a crossing kind is computed per type constructor, with
the element's kind passed through for a list type \citep[\S4.1]{peters2026mode}. OxCaml's
documentation states the corollary the checker relies on: a type crosses uniqueness if it contains no
memory location subject to overwriting \citep{oxcaml2025uniqueness}. For beni we write
$\mathit{crosses}(\tau)$ when $\tau$ is a primitive, or a record, tuple or nominal type all of whose
tracked paths are empty---no \code{List} and no function type is reachable through its fields, through
aliases, or through the payloads of its constructors. A function type never crosses, since its
closure may hold a cell. A type variable does not cross: a value of type \code{a} may be a list, and
it is tracked as a whole, at the single path $\varepsilon$.

\emph{Crossing is decided where the type is declared.} An importer cannot see the constructors of an
opaque type, so it cannot compute crossing itself. The declaring module computes
$\mathit{crosses}(T)$ for each of its types, and publishes it in its interface beside the type, with
the type's parameters' contributions, as Peters et al.\ publish crossing coefficients per type
constructor. For an opaque type that does not cross, the declaring module also publishes its public
functions' summaries over the opaque step $\langle T\rangle$, so an importer sees ``some cell inside
this value'' and never mistakes an opaque value for one that holds nothing. Peters et al.\ report
that without crossing, adoption of modes in a large codebase ``would have been infeasible''
\citep[\S1]{peters2026mode}. Here crossing costs nothing in precision: most values in a beni program
are scalars and records of scalars, and they drop out of the analysis.

\section{Holders, and the two kinds}
\label{sec:model-holders}

A \emph{holder} of a cell $c$ at a program point is a pair $(\mathit{place}, \mathit{path})$
through which $c$ can be reached and read later. Table~\ref{tab:holders} lists the places. The
operand of an update is itself a holder; the check asks about all the others.

\begin{table}[t]
\centering\small
\begin{tabularx}{\linewidth}{@{}>{\raggedright\arraybackslash}p{0.34\linewidth}>{\raggedright\arraybackslash}p{0.2\linewidth}>{\raggedright\arraybackslash}X@{}}
\toprule
Place & Example & Kind \\
\midrule
a local or parameter & \code{xs}; \code{model.rows} as $(\mathit{model}, .\mathit{rows})$ & content \\
a field of a record, tuple or constructor held by a place & $(m, .\mathit{form}.\mathit{tags})$ & content \\
an element of a list held by a place & $(\mathit{rows}, [\ast].\mathit{tags})$ & content \\
an argument already evaluated and not yet passed; a structure whose fields are partly evaluated & \code{xs} in \code{compare xs (List.push xs 1)} & content, until the call or construction that consumes it (\S\ref{sec:model-unique}) \\
the environment of a closure held by a place & $(f, \env)$ & content, live while the closure is \\
the environment of a suspended fiber, a command or subscription thunk the runtime keeps, a \code{Ref} & \code{Cmd.perform (λsend → ⋯ xs ⋯)} & content, \emph{escaped}: no future read can be ruled out \\
a JavaScript sibling, or the result of \code{Js.from} & \code{Js.from xs} & content, escaped \\
a top-level value, with or without a \code{where} clause & \code{empty = []} & content, escaped: read by every later use, for the program's life \\
a value the runtime retains and hands out again: a fiber's outcome (every joiner receives the same value), a bracketed resource (kept for its release), a deferred value's contents & \code{Task.join fib} & content; $\top$ on every path of the result \\
the runtime's hold on a platform-called function's arguments and results & \code{view model} on a runtime that keeps what it rendered & as the platform's entry protocol says (\S\ref{sec:hard-tea}) \\
the runtime's per-row instance item & \code{insts[i].it} & identity only for the runtime (A4); a listener body reads the current item, re-obtained after each dispatch (\Obl{P3}) \\
a derived-value slot & a derived value & content; a list-valued slot is recomputed before it is read (\Obl{P4}) \\
\bottomrule
\end{tabularx}
\caption{The places that can hold a cell.}
\label{tab:holders}
\end{table}

\section{Unique at a site}
\label{sec:model-unique}

The analysis works on an \emph{evaluation-order normal form} of each body, in the style of A-normal
form: every subexpression that does work is bound to a temporary in the order beni evaluates it, and
every call, construction and pattern match takes only variables and literals as operands. A variable
occurrence is a \emph{use at the instruction that consumes it}---the call that receives it as an
argument, the construction that stores it, the match that inspects it---and not at its textual
position. In \code{compare xs (List.push xs 1)} the push is an instruction of its own, evaluated
first; the occurrence of \code{xs} as \code{compare}'s first argument is a use at the \code{compare}
instruction, \emph{after} the push. The same holds for a tuple or record literal whose earlier fields
name a value that a later field edits.

Let $u$ be an \emph{update site}: a call of an update primitive (a \emph{demand},
\S\ref{sec:model-demands}) whose operand is $e$. By beni's evaluation order---the callee, then the
arguments left to right, and the call last---the write happens after every argument of the call has
been evaluated. Let $C(e)$ be the set of cells $e$ may evaluate to.

\begin{definition}[Dead; unique at $u$]
\label{def:unique}
A holder $h$ is \emph{dead} at a program point when no execution continuing from that point uses $h$:
it neither dereferences the cell's elements or length through $h$, nor compares $h$'s reference by
identity, nor passes $h$'s value to anything that may do either. A cell $c \in C(e)$ is \emph{unique
at $u$} when every holder of $c$---the operand's own binding included---is dead at $u$. The operand's
occurrence at $u$ is not a later use: it is consumed by $u$, and the result of the demand is the new
owner of the cell.
\end{definition}

\begin{definition}[The rule]
An update site $u$ is accepted if and only if every $c \in C(e)$ is unique at $u$, $\top \notin
C(e)$ (the checker knows which cells are involved), and no $c$ is \emph{escaped} (held by something
whose future uses the checker cannot see).
\end{definition}

Three consequences follow. First, \emph{a dereference that completes before the write is never
sharing}. \code{List.set xs i (List.get xs j + 1)} reads an element of \code{xs} while evaluating its
third argument, and that read is over before the call; at $u$ the binding \code{xs} is the operand,
and if nothing uses \code{xs} afterwards, the update is accepted. This is Wadler's design goal that
``only update operations are sequentialised but lookups can occur in parallel'' \citep[\S3]{wadler1991use},
and the reservation phase of Rust's two-phase borrows, during which a mutable borrow ``acts exactly
like a shared borrow'' until it is activated, the only new check being at activation \citep{rfc2025}.
What does \emph{not} complete before the write is a value merely held: \code{xs} in
\code{compare xs (List.push xs 1)} is read by \code{compare} after the push, and the program is
rejected---correctly, since under value semantics \code{compare} sees the old list and in place it
would see the new one. Second, a \emph{later} read is sharing: in
\begin{Verbatim}
ys = List.set xs 0 1
List.get xs 0
\end{Verbatim}
the second line reads the old version, and the program is rejected with a message naming the
binding, the \code{set} and the \code{get} (Chapter~\ref{sec:errors}). Third, \emph{a dead alias is
harmless}: \code{saved = xs} followed by \code{List.set xs 0 1}, with \code{saved} never used again,
is accepted---Boyland's condition that aliases need only ``all be dead''
\citep[\S3.2]{boyland2001alias}.

This is not Clean's uniqueness. Clean's \code{*} attribute is a structural property---exactly one
reference exists---checked by counting occurrences: an applied occurrence of a variable is marked
non-unique ``if x either occurs more than once or x is a letrec-variable'', a marking that the
syntax-directed system of Def.~8.14 then uses \citep[\S8]{barendsen1996uniqueness}, refined by the
compiler's observing-reference rule only for guards and strict \code{let}s
\citep[\S9.4]{plasmeijer2011clean}. Liveness accepts reads before a write and dead aliases, which the
marking rejects, with no \code{let!} construct \citep{wadler1990linear}. But the two are
\emph{incomparable}, not ordered: Clean's unique spine with unique elements \citep[\S6]{barendsen1996uniqueness}
types edits of one element while the rest of the list lives, which an abstraction that merges all
positions under $[\ast]$ can accept only with the exclusive-contents flag of \S\ref{sec:model-excl}.

\section{Abstract values}
\label{sec:model-abstract}

For every expression the analysis computes an \emph{abstract value} $A(e)$: a finite map from the
tracked paths under the value to sets of abstract cells,
\[
A(e) : \mathit{Path} \rightharpoonup \wp(\mathit{Cell})
\qquad
\begin{array}{r@{\;}c@{\;}l@{\quad}l}
\mathit{Cell} & ::= & \site(\iota) & \text{an array or a once closure made at instruction } \iota\\
 & \mid & \pi_i.q & \text{the cell at path } q \text{ of parameter } i \text{, on entry}\\
 & \mid & \top & \text{a cell the analysis cannot name}
\end{array}
\]
We write $A(e)@q$ for the set at path $q$. Two conventions keep the domain honest. \emph{$\emptyset$
means provably no cell}---a path of a crossing type, a closure that captured nothing, or a fresh
container known to be empty---and never ``unknown''. \emph{$\top$ is prefix-closed}: a value whose path
$q$ maps to $\top$ has $\top$ at every path under $q$. A decoder's \code{List (List Int)} whose outer
cell is fresh but whose inner cells its declaration does not describe therefore has $[\ast] \mapsto
\{\top\}$, and the result of a call with no summary is $\top$ everywhere. A demand on a path that
holds $\top$ fails the first side condition, \Known\ (\S\ref{sec:model-check}).

The domain is the write-set analysis's abstract domain (\S\ref{sec:bg-browser}) with its identity
facts replaced by cell sets. The analysis is \emph{directional}, in the style of Andersen's
inclusion-based points-to analysis \citep{andersen1994program}: a binding \code{y = x} sets
$A(y) = A(x)$, a copy of the sets, never a merged equivalence class as in unification-based analyses
\citep{steensgaard1996points}. The difference is the precision argument. Emre et al., studying the
translation of C to safe Rust, found that under a unification-based model ``having only four
necessarily unsafe raw pointers is enough to poison 85\,\% of the total raw pointers in the worst
case'' across their benchmarks \citep[\S5]{emre2023aliasing}; on tmux, their largest, adding
directionality alone---going from an equality-based to a subset-based analysis---raised the share of
pointers that affect at most 1\,\% of the others from 7.9\,\% to 88.9\,\%
\citep[\S6, Table~2]{emre2023aliasing}. Hindley--Milner unification is therefore the wrong solver
for this fact; beni's effect bits already live beside the unifier rather than in it, and ownership
does the same.

\section{Transfer rules}
\label{sec:model-transfer}

Table~\ref{tab:transfer} gives the abstract value of each core form of the normal form, with $\Sigma$
the environment of locals. Each rule says what the result may hold and what new holders exist. On
abstract values, $\cup$ is pointwise union, $\mathit{shift}(A, s)$ prefixes every path with step $s$,
and $\mathit{proj}(A, s)$ keeps the paths under $s$ with the step removed. In the table, ``$\cdots$''
is the meta-level ellipsis; beni's own spread is written with three dots.

\begin{table}[t]
\centering\small
\begin{tabularx}{\linewidth}{@{}>{\raggedright\arraybackslash}p{0.25\linewidth}>{\raggedright\arraybackslash}p{0.33\linewidth}>{\raggedright\arraybackslash}X@{}}
\toprule
Form & $A(e)$ & Holders created; notes \\
\midrule
variable $x$ & $\Sigma(x)$ & none new; a \emph{use} at the consuming instruction \\
list literal with elements $x_1 \cdots x_n$ and spreads $s_1 \cdots s_m$ & $\varepsilon \mapsto \{\site(\iota)\} \cup \bigcup_i \mathit{shift}(A(x_i), [\ast]) \cup \bigcup_j \mathit{shift}(\mathit{proj}(A(s_j), [\ast]), [\ast])$ & a fresh cell; a spread's \emph{elements} are stored, its own cell is only read \\
\code{xs ++ ys} on lists & as the literal \code{[ …xs, …ys ]} & a fresh cell; not a demand \\
record, tuple, constructor & $\bigcup_f \mathit{shift}(A(x_f), .f)$ & the structure holds what its parts hold \\
record update \code{\{ r | f = y \}} & $A(r)$ with the paths under $.f$ replaced by $\mathit{shift}(A(y), .f)$ & the other fields are the very same values, so old and new records share every cell under every other field \\
field or tuple access \code{r.f} & $\mathit{proj}(A(r), .f)$ & an alias of the field's cell, not a copy \\
\code{case x of p → e} & per arm, pattern variables bound to projections of $A(x)$; the \code{rest} of \code{[ x, …rest ]} is bound to \emph{the same cells as $x$} (a view of the base, \S\ref{sec:hard-views}); the \code{init} of \code{[ …init, last ]} to a fresh cell & arms are analysed separately; results are joined by union; liveness is per path \\
\code{x = e₁} then \code{e₂} & $\Sigma[x \mapsto A(e_1)]$ & $x$ holds everything in $A(e_1)$ \\
lambda with captures $z_1 \cdots z_m$ & $\env \mapsto \bigcup_j \mathit{flat}(A(z_j))$; and $\varepsilon \mapsto \{\site(\iota)\}$ if the lambda is \emph{once} (\S\ref{sec:model-functions}) & the closure holds every captured cell under $\env$; $\mathit{flat}$ collects all cells of a value, at every path \\
call \code{f a₁ ⋯ aₙ} of a top-level $f$ & by the summary $S_f$ & a use of every argument at the paths the summary reads; a \emph{demand} when some parameter path is consumed \\
call \code{g a₁ ⋯ aₙ} through a value & by $g$'s ownership class (\S\ref{sec:inf-classes}); if $A(g)@\varepsilon$ holds a cell, the call is also a demand on it & as above \\
\code{foreign} or \code{Js} operation with no asserted summary & $\varepsilon \mapsto \{\top\}$, prefix-closed & every cell of every argument is read and becomes escaped (\S\ref{sec:hard-foreign}) \\
top-level value $v$ & $\Sigma(v)$ & every cell escaped: a holder for the program's life (\S\ref{sec:model-check}) \\
\bottomrule
\end{tabularx}
\caption{Transfer rules of the cell analysis.}
\label{tab:transfer}
\end{table}

\section{Function values: environments and affinity}
\label{sec:model-functions}

A closure carries what it captured wherever the closure goes---returned from the function that made
it, stored in a record, a constructor or a list, captured by another closure, passed to a callee. Two
facts about it must travel with it: \emph{which cells it holds}, and \emph{whether it may be called
only once}. A lambda written in the body being analysed shows both; a closure that arrives as the
result of a call, or out of a field, shows neither unless the abstract value carries them. So the
abstract value of every function-typed path carries them, and summaries publish them.

\paragraph{Environments.} The $\env$ step reaches every cell a closure captured. For the lambda in
\begin{Verbatim}
mk : List Int → (⊤ → Int)
mk xs = λ⊤ → List.length xs
\end{Verbatim}
the result's abstract value is $\env \mapsto \{\pi_1\}$, and \code{mk}'s summary says so: its result
path $\env$ aliases its first parameter. A caller that writes \code{g = mk xs} therefore has
$A(g)@\env = A(\mathit{xs})$, and $(g, \env)$ is a holder of \code{xs}'s cell for as long as \code{g}
is live. A placeholder such as \code{List.get xs \_} is a lambda that captures \code{xs}, and is
treated the same way. A closure stored in a record holds its captures under the field's path, so a
record of handler functions holds the lists its handlers captured exactly as long as the record is
live---not for the program's life.

\paragraph{Affinity.} A closure that may \emph{consume} a cell it captured---write it in place, or
pass it to a position that does---must not be called twice, since the second call would see the
first call's write. OxCaml's lock rule says that such a closure is \code{once}
\citep[\S3.3]{lorenzen2024oxidizing}. We make the property one of the closure \emph{value}, not of
the binding that first names it, by representing a once closure as a \emph{cell}: its creation
allocates an abstract cell $\site(\iota)$ at the closure's own path $\varepsilon$, and \emph{every call
of it is a demand on that cell}. Everything the rules already say about cells then applies to once
closures with no new rule: a once closure returned from a function names a fresh cell in the
function's result; one stored in a record is held at the field's path; one captured by another
closure is held under that closure's $\env$; and calling it twice fails \Dead\ at the first call,
because the holder is used again by the second. A closure is once when
\begin{itemize}
\item its body consumes one of its captured cells, or
\item its body calls a captured function value that is once (a demand on a captured cell is a
consumption of it), or
\item its body passes a captured once value to a position that consumes it.
\end{itemize}
Creating a once closure is itself a demand on the captured cells it may consume: the four side
conditions of \S\ref{sec:model-check} are checked at the creation, with the new closure as the only
holder allowed to remain live. From then on the captured cells are reachable only through the
closure, and the closure is used at most once.

A function type in a summary or a class says, as part of its abstract value, whether its values may be
once (a cell at $\varepsilon$) and what they may hold (cells at $\env$). This is mode inference on
arrows---OxCaml's linearity and uniqueness axes in another form---carried by inference and published
in interfaces, never written by the programmer.

\section{Summaries}
\label{sec:model-summary}

The compiler infers a summary for every function and publishes it in the module interface. It is the
annotation the programmer does not write. For a function $f$ with parameters $\pi_1 \ldots \pi_n$ and
result $\rho$:
\[
\begin{array}{r@{\;}l@{\qquad}l}
S_f = \langle & \mathit{use} : (i, q) \mapsto U \subseteq \{\mathsf{R}, \mathsf{A}, \mathsf{E}, \mathsf{W}\} & \text{for each tracked parameter path } \pi_i.q\\
, & \mathit{res} : q \mapsto \wp(\{\pi_i.q'\} \cup \{\fresh_k\}) & \text{for each tracked result path } \rho.q\\
, & \mathit{xres} : q \mapsto \{\mathsf{yes}, \mathsf{no}\} & \text{whether the container at } \rho.q \text{ is exclusive}\\
, & \mathit{xreq} : (i, q) \mapsto \{\mathsf{yes}, \mathsf{no}\} & \text{whether } f \text{ needs } \pi_i.q \text{ exclusive}\\
, & \mathit{fun} : j \mapsto \langle \mathit{calls}, \mathit{supply}, \mathit{need} \rangle & \text{for each function-typed parameter } \pi_j\\
, & \mathit{reason} : (i, q) \mapsto \text{a derivation tag} & \text{why each fact was inferred} \;\rangle
\end{array}
\]
The $\mathit{use}$ of a parameter path is a set of four independent facts:
\begin{itemize}
\item \textsf{R}, \emph{read}: during the call $f$ may dereference the cell.
\item \textsf{A}, \emph{aliased}: the cell may be reachable from $f$'s result, at the paths
$\mathit{res}$ names. The caller accounts for it through $\mathit{res}$.
\item \textsf{E}, \emph{escaped}: the cell may be kept by something that outlives the call and that the
caller cannot see---a command, a fiber, a \code{Ref}, JavaScript, a runtime structure. The caller
marks the cell escaped.
\item \textsf{W}, \emph{written}: $f$ may write the cell in place, directly or by passing it to a
position that does. The caller must prove the cell unique at the call. A written parameter may also
be aliased by the result---every list writer returns its operand---so ownership is transferred, not
lost.
\end{itemize}
We name the common combinations: $\unused = \emptyset$, $\borrowed = \{\mathsf{R}\}$, $\aliased$ when
\textsf{A} is present and \textsf{E}, \textsf{W} are not, $\escaped$ when \textsf{E} is present, and
$\consumed$ when \textsf{W} is present. \emph{Unused} matters: a helper that replaces a field of a
record it receives, \code{changed model todos = \{ model | todos = todos \}}, never reads its first
argument's $.\mathit{todos}$, and a caller's holder passed there is not used by the call. \emph{Read}
is Aspinall, Hofmann and Kone\v{c}n\'y's aspect~3, ``used read-only and not shared with the result'';
\emph{aliased} is their aspect~2, ``used read-only but shared with result''; \emph{written} is their
aspect~1, ``used destructively'' \citep[abstract]{aspinall2008usage}. \emph{Read} is also
FP\textsuperscript{2}'s borrowed parameter, which ``cannot be used in a destructive match, passed as an
owned parameter, or returned as a result'' \citep[\S1.3]{lorenzen2023fp2}. \emph{Escaped} is what
earlier versions of these rules merged with \emph{aliased}; separating them is what lets a caller tell a
cell its result explains from one something invisible keeps.

The result's $\mathit{res}$ records which argument paths each result path may alias. Fresh cells are
\emph{numbered}: two result paths that name the same $\fresh_k$ may hold the same new cell, and two
that name different numbers never do. \code{pair ⊤ = xs = [ 1 ]; \{ a = xs, b = xs \}} has
$\mathit{res}(.a) = \mathit{res}(.b) = \{\fresh_1\}$, so a caller knows the two fields alias. This is
the dependent result qualifier of reachability types, where a function type names the argument its
result reaches, as in \code{id(x: T}$^{\blacklozenge}$\code{): T}$^{\{x\}}$
\citep[\S2.2]{wei2024polymorphic}, with numbered freshness markers. A function-typed result path
carries $\mathit{res}$ at $\varepsilon$ (a once closure's own cell, always fresh) and at $\env$ (what
the closure holds).

$\mathit{fun}(j)$ records, for a function-typed parameter $g$:
\begin{itemize}
\item $\mathit{calls} \in \{0, {\le}1, \mathit{many}\}$: how often $f$ may call $g$ (a call inside a
loop, a recursion or a callee that calls its own callback many times is $\mathit{many}$). A once $g$ is
accepted only where $\mathit{calls} \le 1$; this is the summary-level form of the rule that a second
call of a once closure fails \Dead.
\item $\mathit{supply}(k)$: what $f$ passes as $g$'s $k$th argument, as a set of \emph{provenances}---paths
$\pi_i.q$ of $f$'s own parameters, or $\fresh$ for a cell $f$ made---together with an ownership bit,
\emph{owned} when the passed cell is unique at the call of $g$ and \emph{borrowed} otherwise. The
provenance is what lets $g$'s own $\mathit{res}$, which names $g$'s parameters, be translated into
$f$'s parameters at a call of $f$ (\S\ref{sec:inf-classes}).
\item $\mathit{need}$: what $f$ requires of $g$'s result---nothing, or that it be fresh and exclusive, when
$f$ stores it where exclusivity is promised.
\end{itemize}
The ${\le}1$/$\mathit{many}$ distinction is OxCaml's \code{once}/\code{many}: when a unique value is
consumed in a closure, ``this closure can be invoked at most once''
\citep{oxcaml2025uniqueness,lorenzen2024oxidizing}. It is also Futhark's reason for refusing to let a
lambda consume a variable bound outside it: the compiler has no idea how often the function will be
applied \citep{henriksen2022uniqueness}. An earlier version of these rules also kept an
$\mathit{escapes}$ bit per function-typed parameter; it is now the $\mathsf{E}$ fact on the
parameter's own $\env$ and $\varepsilon$ paths.

\paragraph{Conditional summaries.} Summaries may depend on the facts of a function-typed parameter's
class. \code{List.foldl}'s accumulator is written only when its callback writes its own accumulator
parameter: $\mathsf{W} \in \mathit{use}(\mathit{acc})$ iff $\mathsf{W} \in \mathit{use}_g(2)$. A body
whose check depends on whether a closure handed to a parameter is kept, such as
\begin{Verbatim}
f xs h =
    g = λ_ → List.length xs
    _ = h g
    List.push xs 1
\end{Verbatim}
is accepted iff $\mathsf{E} \notin \mathit{use}_h(1)@\env$, and that condition travels in $f$'s
summary to be checked at each caller when $h$'s class is known. Formally, every fact of a summary and
every side condition may be a \emph{guarded} value: a join $\bigsqcup_g (\gamma_g \Rightarrow v_g)$ in
which each guard $\gamma_g$ is a conjunction of \emph{atoms}, an atom being one fact of one
function-typed parameter's class (``$\mathsf{W} \in \mathit{use}_h(1)$'', ``$h$ may be once''). These are
rank-1 constraints, the same shape the effect summaries already have, in which a class may depend on
``the other classes of the same scheme that reach it''. Their lattice and its height are given in
\S\ref{sec:inf-lattice}.

\paragraph{The bracket notation.} We write summaries on a type, in a notation of this thesis only:
\begin{center}\small
\begin{tabular}{@{}r@{\;}c@{\;}l@{}}
\textit{row} & ::= & \textit{param}, $\cdots$, \textit{param} \code{→} \textit{result}\\
\textit{param} & ::= & \textit{type} \code{[} \textit{fact} \code{;} $\cdots$ \code{]}\\
\textit{fact} & ::= & \textit{use} $\mid$ \code{calls} $n$ $\mid$ \code{supply} $k$ \textit{own} \textit{prov} $\mid$ \code{need fresh excl} $\mid$ \code{requires excl} $\mid$ \textit{fact} \code{if} \textit{guard}\\
\textit{use} & ::= & \code{unused} $\mid$ \code{borrowed} $\mid$ \code{aliased} $\mid$ \code{escaped} $\mid$ \code{consumed} $\mid$ \textit{use}\code{+}\textit{use}\\
\textit{result} & ::= & \textit{type} \code{[=} \textit{prov}, $\cdots$ \code{]} $\mid$ \textit{type} \code{[=} \textit{prov}, $\cdots$ \code{; excl]}\\
\textit{prov} & ::= & \code{π}$_i$\textit{path} $\mid$ \code{fresh} $\mid$ \code{fresh}$_k$ $\mid$ $g$\code{.res}\\
\textit{own} & ::= & \code{owned} $\mid$ \code{borrowed}\\
\textit{guard} & ::= & \textit{atom} $\mid$ \textit{guard} \code{and} \textit{guard}
\end{tabular}
\end{center}
A parameter with no bracket is $\unused$ or crosses; a result with no bracket is crossing.

\section{Demands: declared, not inferred}
\label{sec:model-demands}

\paragraph{The writers, and why syntax is not one.} The demand set is the core library's \emph{named
writers}: \code{set}, \code{update}, \code{push}, \code{pop}, \code{swap}, \code{insertAt},
\code{removeAt}, \code{cons} and \code{append}, each called by name. List \emph{syntax} is not a demand. A
literal \code{[ x, …xs ]}, a concatenation \code{[ …xs, …ys ]} and \code{xs ++ ys} on lists build a
new array and only read their operands, as the same expressions do in JavaScript. An earlier version
of these rules lowered spreads and \code{++} to the consuming \code{cons} and \code{append}. That made
every spread of a list the program still used an error: a top-level list could never be spread, a
\code{view} could not concatenate two model lists, and \code{a = base ++ [ 1 ]; b = base ++ [ 2 ]}
was rejected at the first line. Those are among the most frequent list operations in Elm-style code.
A construction that copies is visible in the source and its cost is the cost a hand-written program
pays; it is not the \emph{silent} copy that this design refuses, which replaces an in-place write the
programmer asked for. Where the operand of a spread or \code{++} is unique and dead at the site---so
that the named writer would have been accepted---the checker emits a \emph{warning} naming the writer
(\S\ref{sec:disc-rejected}); it never changes what the program does.

\paragraph{Assertions.} The list writers in beni's core library are written in beni over the
low-level \code{Js} interface (\S\ref{sec:bg-boundary}): \code{set xs i v}, for instance, is a
\code{Js.pure} computation that slices the base array through \code{Js.call} and writes the copy. By
the transfer rules, a value passed to \code{Js} is escaped and a \code{Js} result is $\top$, so
inference would conclude that every writer ``escapes its operand and returns an unknown cell''. The
writers therefore carry an \emph{asserted} summary, a promise the body cannot show, just as a
\code{foreign} declaration carries an asserted effect rung and its \code{sync} marks. The core library
and the platforms may assert a summary; every other package may only \emph{narrow} an inferred one
(\S\ref{sec:disc-options}, Change~5). A false assertion is a correctness bug in privileged code, not a
user error. The same applies to every function in core or a platform whose body goes through
\code{Js} or a \code{foreign} and handles a cell, not only to the writers;
\S\ref{sec:soundness-trusted} counts that surface. The writers' rows are:
\begin{Verbatim}
set      : List a [consumed+aliased], Int, a [aliased]
             → List a [= π₁; [*] = π₁[*], π₃; excl if π₁ excl and π₃ stored exclusively]
update   : List a [consumed+aliased; requires excl if g consumes], Int,
           (a → a) [calls ≤1; supply 1 owned π₁[*] if π₁ excl; supply 1 borrowed π₁[*] otherwise]
             → List a [= π₁; [*] = π₁[*], g.res; excl if π₁ excl and g.res fresh]
push     : List a [consumed+aliased], a [aliased]
             → List a [= π₁; [*] = π₁[*], π₂; excl if π₁ excl and π₂ stored exclusively]
pop      : List a [consumed+aliased] → List a [= π₁; excl if π₁ excl]
swap     : List a [consumed+aliased], Int, Int → List a [= π₁; excl if π₁ excl]
insertAt : List a [consumed+aliased], Int, a [aliased]
             → List a [= π₁; [*] = π₁[*], π₃; excl if π₁ excl and π₃ stored exclusively]
removeAt : List a [consumed+aliased], Int → List a [= π₁; excl if π₁ excl]
cons     : a [aliased], List a [consumed+aliased]
             → List a [= π₂; [*] = π₂[*], π₁; excl if π₂ excl and π₁ stored exclusively]
append   : List a [consumed+aliased], List a [borrowed; [*] aliased]
             → List a [= π₁; [*] = π₁[*], π₂[*]]
length   : List a [borrowed] → Int
copy     : List a [borrowed; [*] aliased] → List a [= fresh; [*] = π₁[*]]   -- new
at       : List a [borrowed; [*] aliased], Int → a [= π₁[*]]               -- core-private
\end{Verbatim}
A stored element is aliased by the result at $[\ast]$, and ``stored exclusively'' means unique at the
store and disjoint from every cell already in the list, as \S\ref{sec:model-excl} defines. The $\excl$
annotations are the subject of that section; \code{append} and \code{copy} never keep the flag,
because their results share elements with a list that lives on. An index is an \code{Int}, which crosses and so carries no bracket. \code{update} is
\code{set} with a callback, and its callback is supplied an \emph{owned} element only when the list's
contents are exclusive. Everything else in the list library is ordinary beni over \code{at}, the array
builder and a few internal primitives with asserted rows, and its summary is inferred: \code{map}, for
instance, comes out as
\begin{Verbatim}
map : List a [borrowed], (a → b) [calls many; supply 1 borrowed π₁[*]] → List b [= fresh; [*] = g.res]
\end{Verbatim}
so that \code{List.map model.rows (λr → r.tags)} is accepted and its result's elements are known to
alias \code{model.rows}'s elements' \code{tags}. Putting the demand only on the named writers follows a
Clean recipe: library functions such as \code{newArray} return an array whose attribute is a variable
rather than ``unique'', because ``if we are careful never to return a unique array from a function, we
will always be able to share arrays'' \citep[\S6]{devries2008simplified}. User code carries nothing.

\paragraph{A summary describes the emitted code, and every alias the library returns is in it.}
beni's backend compiles a function that builds its result in tail position modulo a list constructor,
such as
\begin{Verbatim}
appendTo xs ys = case xs of
    [] → ys
    [ x, …rest ] → [ x, …appendTo rest ys ]
\end{Verbatim}
into a builder loop whose exit appends \code{ys}'s elements to the builder---and today returns
\code{ys} \emph{itself} when the builder is empty (\S\ref{sec:bg-lists}). The checker reads such a
function as the backend emits it. Under the rule the exit is changed to \emph{copy} \code{ys} when the
builder is empty, so that the result is always fresh; the copy costs $O(|\mathit{ys}|)$, which is what
the non-empty case already pays to append \code{ys}'s elements, so the loop's cost does not change. The
alternative, a result $\{\fresh, \pi_2\}$, would make every such function's result alias its last
argument, and is worse. The same audit applies to every function in the library that may return an
input or a view of it (\S\ref{sec:bg-lists}): each either says so in $\mathit{res}$ or stops doing it.
We propose (\S\ref{sec:disc-options}, Change~2):
\begin{itemize}
\item functions that build a new list always return a fresh array: \code{filter}, \code{map},
\code{slice}, \code{take}, \code{reverse}, the sorts, \code{concat}, and the builder exit above;
\item functions that return a view keep doing so and say $\mathit{res} = \{\pi_1\}$: \code{drop} and
\code{tail}, like a pattern's \code{rest};
\item the writers' no-ops (an out-of-range \code{set}, \code{swap xs i i}, \code{pop []}) already return
$\pi_1$, which is what their rows say;
\item the claimed-head prepend disappears with the trie, and \code{cons} on a unique array writes it.
\end{itemize}
Loops whose list traversal the backend rewrites into an offset into the base array read the base, and
are summarised as such reads.

\section{Liveness, per path}
\label{sec:model-liveness}

$\live(h, p)$ holds for a holder $h = (x, q)$ at point $p$ when there is a use of $x$ at or after $p$,
on some path, that reads $q$ or a prefix of it. The uses, in evaluation order, are:
\begin{itemize}
\item a read of $x.q'$ with $q' \sqsubseteq q$ or $q \sqsubseteq q'$ (prefix order);
\item $x$ passed to a call whose summary's $\mathit{use}$ at a path that meets $q$ is not $\unused$; a
parameter with no summary reads everything;
\item $x$ stored in a structure that is live, returned, or captured by a closure that is live---a
captured variable is live exactly while the closure's $\env$ is, and forever once the closure
escapes;
\item a call of $x$ when $x$ is a function value: a use of everything under $(x, \env)$;
\item $x$ used whole by anything the analysis does not see into: a \code{Js} operation or
\code{foreign} with no asserted summary is a read of every path, and its arguments' cells become
escaped (Table~\ref{tab:transfer} says the same);
\item \code{==} and \code{<}, which are method calls resolved by the dispatch table: a type's \code{eq}
and \code{compare} only read their arguments, by the rule of \S\ref{sec:inf-generic}, so the call
reads every path of both operands and writes none.
\end{itemize}
A \code{case} makes liveness path-sensitive in the usual way: a holder live only in the arm not
taken is dead in the other (\S\ref{sec:hard-branches}). A self tail call rebinds the parameters: each
new argument reads the parameters as they were, and the old value of a parameter is dead after the
jump unless an argument still to be evaluated uses it. The backend's freedom to rebind a parameter
early applies only when no argument still to come reads it, so it never moves a use across a rebinding.

Liveness is computed per (variable, path), so that \code{m.name} read after
\code{List.push m.rows x} is not a use of \code{m.rows}. Field-sensitive deadness is what lets a
record be threaded through helpers that each touch one field. In our prototype study its absence
cost a copy of every row on each append (\S\ref{sec:bg-measure}).

\section{The check}
\label{sec:model-check}

At a call $f\,\bar a$ whose summary has $\mathsf{W} \in \mathit{use}(i, q)$, and at a call of a
function value $g$ with a cell $c \in A(g)@\varepsilon$, for each such cell $c$:

\begin{description}[leftmargin=1.2em,font=\normalfont]
\item[\Known] $c \neq \top$: the cell is known.
\item[\Unescaped] $c$ is not escaped: nothing whose uses the checker cannot see holds it.
\item[\Dead] for every holder $h$ of $c$, $\neg\live(h, \text{after the call})$: nothing uses it later.
\item[\Disjoint] for every other argument path $(a_j, q')$ with $j \neq i$ or $q' \neq q$---including
other paths of the \emph{same} argument---and for every cell under $\env$ of every function-valued
argument, $c \notin A(a_j)@q'$: nothing reaches it \emph{during} the call.
\end{description}

\Dead\ reads ``after the call'': every argument has been evaluated, so a dereference inside an argument
is complete before the write; a variable merely passed to a later call is a use at that later call.
The operand's own binding is a holder like any other, so \code{ys = List.set xs 0 1} followed by
\code{List.get xs 0} fails \Dead\ on \code{xs}. \Disjoint\ is the separation rule. The write happens
\emph{inside} the callee, before the callee's own later reads, and the callee reads its other
parameters and calls its callbacks while the write is already visible. In
\code{List.foldl xs xs (λx acc → List.insertAt acc 0 x)} the walked input and the consumed accumulator
are the same array. In \code{g xs (λ\_ → List.length xs)} with \code{g xs k = k (List.push xs 1)}, the
callback would read the pushed array through its capture. In \code{f r} where \code{r.a} and \code{r.b}
are one cell and \code{f} pushes \code{π₁.a} and then reads \code{π₁.b}, the two paths of one argument
alias. \Disjoint\ corresponds to Rust's rule of one mutable or any number of shared borrows at a call
\citep{rfc2094}, to Futhark's requirement that the argument for a consumed parameter not alias any
other argument \citep{henriksen2017thesis}, and to Aspinall, Hofmann and Kone\v{c}n\'y's reading of the
comma in a typing context as a tensor: two components ``must be implemented without sharing, unless
the access is guaranteed to be read-only'' \citep[\S3]{aspinall2008usage}. It is checked on abstract
cells, so two arguments that \emph{may} be one cell fail it.

When the operand is itself a parameter path $\pi_j.q'$ of the enclosing function, that function's
$\mathit{use}(j, q')$ gains \textsf{W}, and \Unescaped, \Dead\ and \Disjoint\ for $\pi_j.q'$ are
discharged at \emph{its} callers in turn---including \Disjoint\ against the caller's other paths of
the same argument, which is how the function may treat its parameter paths as distinct cells.
A failure of \Known\ is reported as a \emph{limit} of the checker; a failure of \Unescaped\ as a limit
or as sharing, according to the tag of the escape (Table~\ref{tab:tags}); a failure of \Dead\ or
\Disjoint\ as sharing, with the holder named (Chapter~\ref{sec:errors}).

\paragraph{Holders that live for the whole program.} A top-level value, including one with a
\code{where} clause that is memoised per evidence, is a holder that is never dead: its cells are
escaped from the start, so \code{List.push defaults x} on a top-level list is rejected, with
\code{List.copy} as the fix. Spreading it, \code{[ x, …defaults ]}, only reads it and is accepted. The
runtime's hold on the model between dispatches is the other such holder; the entry protocol of
\S\ref{sec:hard-tea} says when it ends.

\section{Deep uniqueness: exclusive contents}
\label{sec:model-excl}

A list's elements and a record's fields can hold cells, and two positions of one container can hold
the \emph{same} cell: \code{[ r, r ]}, \code{List.repeat r 3}, a dictionary into which one list was
inserted under two keys, or one element whose two fields hold one list,
\code{[ \{ a = xs, b = xs \} ]}. An edit inside an element would then be seen at the other position.
The cell-set domain cannot tell positions apart: every element of a list built in a loop comes from
one allocation site, so $A(\mathit{rows})@[\ast].\mathit{tags}$ is one abstract cell for all of them,
and \Dead\ finds the holder $(\mathit{rows}, [\ast].\mathit{tags})$ live---the other elements are read
later---whenever the container outlives the edit. That is the right answer when positions may share
and a false error when they cannot. It is the case Clean handles with \emph{deep} uniqueness, a unique
list of unique elements, from which an element may be handed out unique because the list's
construction guaranteed that no element was shared when it went in \citep[\S9.2]{plasmeijer2011clean};
Barendsen and Smetsers enumerate the admissible attributions of a list constructor for this purpose,
the third being ``for lists of which both the spine and the elements are unique''
\citep[\S6]{barendsen1996uniqueness}. OxCaml likewise treats uniqueness as ``a deep property by
default; that is, we expect components of a unique value also to be unique''
\citep[\S2.1]{lorenzen2024oxidizing}.

We add a flag $\excl$ to abstract container paths: $(x, q[\ast])$ for a list, and a folded recursive
path such as a dictionary's values. A container path is \emph{exclusive} at a point when:
\begin{enumerate}
\item \emph{across positions}: no cell is reachable from two distinct positions of the container;
\item \emph{within a position}: within each element, every two distinct tracked paths hold distinct
cells, recursively through nested containers;
\item \emph{no outside holder}: no cell reachable through the container has a live holder other than
the container's own path.
\end{enumerate}
Condition 2 is the one an earlier version of these rules left out; without it, the element
\code{\{ a = xs, b = xs \}} counted as exclusive, and an edit of \code{r.a} inside \code{List.update}'s
callback was seen through \code{r.b}. It is what an ordinary call gets from the caller's \Disjoint\
across paths of one argument, and an owned supply from a container must give the callback the same
guarantee. The flag is
\begin{itemize}
\item \emph{set} by a construction every stored value of which is unique at the store \emph{and}
pairwise disjoint from every other stored path, checked as \Disjoint\ over all tracked paths of all
stored elements: a literal of distinct fresh values; a fresh result whose summary says $\excl$
(\code{map}'s result when its callback's result is fresh at every tracked path---each call allocates
its own cells---and distinct across paths; a decoder's output, by its declaration);
\item \emph{kept} by a \emph{destructive read}: a container function that consumes the container,
hands one element to a callback and stores the callback's result back at that position, reading no
other element in between. \code{List.update}'s body is \code{set xs i (func (at xs i))}, and
\code{set} writes slot $i$. The element is unique inside the callback because the container is
consumed at the call (\Dead), no argument or capture of the callback holds it (\Disjoint), and by
conditions 1--3 its own paths are disjoint and held only through the slot, which nothing reads before
it is overwritten. The flag survives if the callback's result is fresh or the element itself. A
dictionary's \code{update}, which does the same at a \emph{key}, has the same shape; since the path
domain cannot name a key, its callback's $\mathit{supply} = \mathrm{owned}$ is an asserted row with the
body-level reason written beside it;
\item \emph{kept or cleared by every writer} as its row in \S\ref{sec:model-demands} says: a store of a
unique value disjoint from the rest keeps it, a store of anything else clears it, and \code{append} and
\code{copy} clear it because their results share elements with a list that lives on;
\item \emph{split} by a pattern \code{[ x, …rest ]} on an exclusive container whose scrutinee is dead
after the match: \code{x}'s cells and the cells under \code{rest}'s $[\ast]$ are then known to be
disjoint, and \code{rest} stays exclusive. Without the split, \code{x} and \code{rest} are aliases at
$[\ast]$, because they come from one abstract position;
\item \emph{cleared} by a non-destructive read whose result is live while the container is (a dictionary
lookup whose dictionary is used afterwards, \code{List.get}, the \code{x} of a pattern while the
scrutinee lives), and by any join with a path that is not exclusive.
\end{itemize}
A function that needs a parameter's contents exclusive---because it hands an element to a writer, or
splits and edits the head---records $\mathit{xreq}(i, q) = \mathsf{yes}$, and every caller checks the
flag at the call, as it checks uniqueness. So
\begin{Verbatim}
bumpHead : List (List Int) → List (List Int)
bumpHead rows = case rows of
    [] → []
    [ r, …rest ] → [ List.push r 0, …rest ]
\end{Verbatim}
is accepted, with $\mathit{use}(\pi_1.[\ast]) \ni \mathsf{W}$ and $\mathit{xreq}(\pi_1) = \mathsf{yes}$,
and each caller must pass a list whose elements are distinct, unshared cells; a caller that passes
\code{[ r, r ]} is rejected, correctly, since the push would be seen through \code{rest}. Clean, with a
unique spine and unique elements, types this program; with the flag, so does the checker.

With the flag, \code{List.update rows i (λr → \{ r | tags = List.push r.tags t \})} is accepted when
\code{rows} is unique and exclusive, and rejected---``the elements of \code{rows} may share''---when they
came from somewhere the checker cannot see. Without the flag the whole class of nested edits would be a
limit. The flag travels in summaries ($\mathit{xres}$, $\mathit{xreq}$) and in interfaces, and its
soundness is Lemma~\ref{lem:excl}. Elements of a crossing type need none of this.
